\documentclass[11pt]{article}
\usepackage[margin=1in]{geometry}
\usepackage{amsmath, amssymb}
\usepackage{listings}
\usepackage{xcolor}
\usepackage{hyperref}

\lstset{
  basicstyle=\ttfamily\small,
  backgroundcolor=\color{gray!5},
  frame=single,
  framerule=0.4pt,
  rulecolor=\color{gray!50},
  breaklines=true,
  columns=fullflexible,
  keepspaces=true,
  showstringspaces=false,
  keywordstyle=\color{blue!70!black}\bfseries,
  commentstyle=\color{green!50!black},
  stringstyle=\color{red!60!black},
  morekeywords={FUNCTION,END,DECLARE,AS,RETURN,WHILE,TRUE,FOR,IF,ELSE,END FUNCTION,END WHILE,END FOR,END IF},
  literate={←}{$\leftarrow$}2 {→}{$\rightarrow$}2
}

\title{SDR Monitoring --- Pseudocode Example}
\author{pseudocode-specific skill}
\date{\today}

\begin{document}
\maketitle

\section*{Overview}

High-level SDR monitoring pipeline: SoapySDR $\rightarrow$ IQ acquisition $\rightarrow$ Welch PSD $\rightarrow$ Plot.

\begin{lstlisting}[language={},caption={High-level flow (pseudocode-general)}]
1. IMPORT LIBRARIES (SoapySDR, numpy, scipy.signal, matplotlib)
2. sdr = SETUP SDR DEVICE (driver="rtlsdr", sample_rate=2.4e6, center_freq=100e6, gain=40)
3. iq = ACQUIRE IQ STREAM (sdr, num_samples=65536)
4. f, psd = APPLY WELCH METHOD (iq, fs=2.4e6, nperseg=1024, noverlap=512)
5. PLOT frequency spectrum (f, psd, xlabel="Frequency (Hz)", ylabel="PSD (dB/Hz)")
6. SAVE figure ("output/sdr_welch.png")
7. LOOP forever:
   7.1 iq = ACQUIRE IQ STREAM (sdr, num_samples=65536)
   7.2 f, psd = APPLY WELCH METHOD (iq, fs=2.4e6, nperseg=1024, noverlap=512)
   7.3 UPDATE plot (f, psd)
   7.4 PAUSE (0.1s)
\end{lstlisting}

\newpage

\section*{Detailed Pseudocode}

\begin{lstlisting}[caption={SDR Monitoring --- pseudocode-specific conventions}]
DECLARE SAMPLE_RATE AS FLOAT <- 2.4e6
DECLARE CENTER_FREQ  AS FLOAT <- 100e6
DECLARE GAIN         AS FLOAT <- 40
DECLARE NUM_SAMPLES  AS INTEGER <- 65536
DECLARE NPERSEG      AS INTEGER <- 1024
DECLARE NOVERLAP     AS INTEGER <- 512

FUNCTION SetupSDRDevice(driver, sample_rate, center_freq, gain)
    DECLARE sdr AS DEVICE
    <- SoapySDR.Device(driver)
    <- sdr.setSampleRate(sample_rate)
    <- sdr.setFrequency(center_freq)
    <- sdr.setGain(gain)
    RETURN sdr
END FUNCTION

FUNCTION AcquireIQStream(sdr, num_samples)
    DECLARE iq     AS ARRAY
    DECLARE buffer AS ARRAY of COMPLEX
    <- sdr.setupStream(SOAPY_RX)
    <- sdr.readStream(buffer, num_samples)
    <- buffer.astype(COMPLEX64)
    RETURN iq
END FUNCTION

FUNCTION ApplyWelchMethod(iq, fs, nperseg, noverlap)
    DECLARE f   AS ARRAY
    DECLARE psd AS ARRAY
    (f, psd) <- scipy.signal.welch(iq, fs=fs,
                 nperseg=nperseg, noverlap=noverlap)
    RETURN (f, psd)
END FUNCTION

FUNCTION PlotSpectrum(f, psd, title)
    DECLARE fig AS FIGURE
    <- matplotlib.pyplot.figure()
    <- matplotlib.pyplot.plot(f, psd)
    <- matplotlib.pyplot.xlabel("Frequency (Hz)")
    <- matplotlib.pyplot.ylabel("PSD (dB/Hz)")
    <- matplotlib.pyplot.title(title)
    <- matplotlib.pyplot.grid(TRUE)
    RETURN fig
END FUNCTION

FUNCTION SaveFigure(fig, path)
    <- fig.savefig(path)
END FUNCTION

FUNCTION ContinuousMonitoring(sdr)
    DECLARE f   AS ARRAY
    DECLARE psd AS ARRAY
    DECLARE fig AS FIGURE
    WHILE TRUE = TRUE
        iq <- AcquireIQStream(sdr, NUM_SAMPLES)
        (f, psd) <- ApplyWelchMethod(iq, SAMPLE_RATE,
                      NPERSEG, NOVERLAP)
        fig <- PlotSpectrum(f, psd, "SDR Monitoring")
        matplotlib.pyplot.pause(0.1)
    END WHILE
END FUNCTION

--- Main ---
sdr <- SetupSDRDevice("rtlsdr", SAMPLE_RATE, CENTER_FREQ, GAIN)
iq <- AcquireIQStream(sdr, NUM_SAMPLES)
(f, psd) <- ApplyWelchMethod(iq, SAMPLE_RATE, NPERSEG, NOVERLAP)
fig <- PlotSpectrum(f, psd, "SDR Monitoring")
SaveFigure(fig, "output/sdr_welch.png")
ContinuousMonitoring(sdr)
\end{lstlisting}

\end{document}
